\chapter{Results, Validation and Conclusion}

\section{Archived device measurements}

The archived suite reports results on a Galaxy S23 with Snapdragon 8 Gen 2,
8 GB of RAM, Android 16 and five XNNPACK threads. The accompanying
presentation reports 16 steps. Build identifiers, exact prompts and
thermal logs for individual cases were unavailable. The suite includes
English, Hindi, Telugu and Tamil. Table~\ref{tab:measurements} summarizes the
cases, with RTF carried forward as reported. The tests used short, medium and paragraph prompts.

\begin{table}[htbp]
\centering
\caption{Archived Galaxy S23 results by language and prompt length}
\label{tab:measurements}
\small
\begin{tabular}{@{}llrrrr@{}}
\toprule
Prompt & Language & Audio (s) & Generation (s) & RTF & Peak RSS (MB) \\
\midrule
Short & English & 2.48 & 3.82 & 1.54 & 378 \\
Short & Hindi & 2.31 & 3.79 & 1.64 & 379 \\
Short & Telugu & 2.52 & 3.85 & 1.52 & 381 \\
Short & Tamil & 2.40 & 3.81 & 1.58 & 380 \\
\addlinespace
Medium & English & 7.12 & 10.84 & 1.52 & 384 \\
Medium & Hindi & 6.85 & 10.61 & 1.54 & 385 \\
Medium & Telugu & 7.30 & 11.02 & 1.50 & 386 \\
Medium & Tamil & 7.05 & 10.79 & 1.53 & 385 \\
\addlinespace
Paragraph & English & 15.82 & 23.85 & 1.50 & 387 \\
Paragraph & Hindi & 16.20 & 24.62 & 1.51 & 388 \\
Paragraph & Telugu & 15.65 & 23.78 & 1.51 & 389 \\
Paragraph & Tamil & 16.10 & 24.47 & 1.52 & 388 \\
\bottomrule
\end{tabular}
\end{table}

Across the 12 archived cases, RTF ranged from 1.50 to 1.64 and peak
resident memory from 378 to 389 MB. Later screenshots show RTF of 7.88
and 11.77 and displayed memory of 933 and 812 MB, respectively. These
screenshots show readings from saved clips with unknown build and device
conditions. They document separate runs with different measurement scope.

\section{Interpret current app telemetry}

Real-time factor (RTF) is compute time divided by output audio duration.
\texttt{MainActivity.kt} divides the sum of native chunk times by the
final stitched waveform duration. An RTF above one means generation took
longer than playback under that timing definition. Reference preparation,
Kotlin stitching and WAV writing fall outside the native chunk timer.
Full latency is the wall-clock interval from tapping Generate to save
completion.

\texttt{TelemetryMonitor.kt} reads \texttt{VmRSS} after stitching and before
saving. This reading is a point sample of resident memory. Peak RSS
requires sampling throughout the run. Model file size measures disk storage
and is reported separately.

\section{Tests and practical checks}

\texttt{MultilingualChunkerTest.kt} checks segmentation, punctuation and
abbreviations. \texttt{WaveformStitcherTest.kt} checks assembly, trimming
and normalization. Device verification should include a short sentence
and a paragraph in airplane mode after installation, followed by playback,
export and sharing. Listening to the resulting WAV can reveal missing
words and clipped endings.

To compare builds, record model hashes, device and Android version, input
audio and texts, sampling steps, threads and whether the run started cold.
Capture native timing, end-to-end time and peak memory separately. The
measurements in this report were carried forward from existing results.

\section{Conclusion and future work}

The project adapted pretrained IndicF5 and Vocos to a fixed-window LiteRT
deployment on Android. Chunking, cached reference features and waveform
stitching let the app generate paragraphs offline and save them as one
WAV. Future work should reproduce latency and peak memory on the latest
build, test acceleration and lower-step sampling, and evaluate naturalness,
intelligibility and speaker similarity.
